\documentclass[svgnames]{article}    

\usepackage{graphicx}                 

% -- math                               
\usepackage{amssymb}
\usepackage{amsmath}
\usepackage{esint}
\usepackage{geometry}

% -- noindent
\setlength\parindent{0pt}

% -- images
%\graphicspath{{ }}                   

% -- tikz
\usepackage{pgfplots}
\pgfplotsset{compat=1.15}
\usepackage[most]{tcolorbox}

\newtcolorbox{mybox}{
    enhanced,
    boxrule=0pt,frame hidden,
    colback=green!10!white,
    sharp corners
}

% -- figures
\usepackage{float}
\usepackage{caption}
\usepackage{lipsum}
\usepackage{comment}

% -- code 
\usepackage{listings}
\usepackage{xcolor} 

\definecolor{codegreen}{rgb}{0,0.6,0}
\definecolor{codegray}{rgb}{0.5,0.5,0.5}
\definecolor{codepurple}{rgb}{0.58,0,0.82}
\definecolor{backcolour}{rgb}{0.95,0.95,0.92}

\lstdefinestyle{mystyle}{
    backgroundcolor=\color{backcolour},   
    commentstyle=\color{codegreen},
    keywordstyle=\color{magenta},
    numberstyle=\tiny\color{codegray},
    stringstyle=\color{codepurple},
    basicstyle=\ttfamily\footnotesize,
    breakatwhitespace=false,         
    breaklines=true,                 
    captionpos=b,                    
    keepspaces=true,                 
    numbers=left,                    
    numbersep=5pt,                  
    showspaces=false,                
    showstringspaces=false,
    showtabs=false,                  
    tabsize=2
}

\lstset{style=mystyle}

% -- start 
\title{Programming 4H-Silicon Carbide's Hamiltonian}
\author{Deval Deliwala}
%\date{}                              

\begin{document}
\maketitle
\tableofcontents 

\vfill
In this paper I define the basis of $\uparrow, \downarrow$'s as the
\textit{Zeeman Basis}. I define the basis with the electron's coupled and the
nucleus's still in the Zeeman basis as the \textit{coupled basis}. I.e. 

\[
|\uparrow \uparrow\rangle|\uparrow \uparrow\rangle \quad \Longleftrightarrow \quad
|11\rangle| \uparrow \uparrow\rangle 
\] \vspace{3px}

defines a conversion between the \textit{Zeeman} and \textit{Coupled} bases.\\ 

See section \ref{Basis} for a complete description of both bases.  
\vfill


\newpage
\section{Zeeman Hamiltonian}

The Zeeman Hamiltonian is defined in the coupled-basis via 

\[
    \hat{H}_Z |s,m\rangle|m_{I_1}\rangle |m_{I_2}\rangle = (mg_e\mu_BB_0 +
    m_{I_1}g_{n_1}\mu_NB_0 + m_{I_2}g\mu_NB_0 )|s, m\rangle |m_{I_1}\rangle
    |m_{I_2}\rangle. 
\] \vspace{3px}

Hence, the coupled-basis is itself the eigenbasis. To build the Hamiltonian, we
first define the coupled basis in the following manner:

\begin{lstlisting}[language=Python]
import sympy as sp

""" build coupled-basis """ 

# electron |s, m> pairs 
electron_states = [
    (1, +1),
    (1,  0),
    (0,  0),
    (1, -1),
]

# nuclear |m_{I_1}>|m_{I_2}> pairs, s=1/2 implicit 
nuclear_pairs = [
    ( sp.Rational(+1,2), sp.Rational(+1,2) ), # (+1/2, +1/2) 
    ( sp.Rational(+1,2), sp.Rational(-1,2) ), # (+1/2, -1/2) 
    ( sp.Rational(-1,2), sp.Rational(+1,2) ), # (-1/2, +1/2) 
    ( sp.Rational(-1,2), sp.Rational(-1,2) ), # (-1/2, -1/2)
]
\end{lstlisting}

Finally, we apply $\hat{H}_Z$ to every basis state and generate the 16$\times$16 
Zeeman Hamiltonian via \verb|sp.diag| since the coupled-basis is already the
eigen (diagonal)-basis. 

\begin{lstlisting}[language=Python] 
# zeeman frequencies
omega_e   = g_e  * mu_B * B0
omega_n1  = g_n1 * mu_N * B0
omega_n2  = g_n2 * mu_N * B0

# build the diagonal entries
entries = []
for s, m in electron_states: 
    for mI1, mI2 in nuclear_pairs:
        entries.append(
            m*omega_e     +      # electron zeeman
            mI1*omega_n1  +      # nucleus 1 zeeman
            mI2*omega_n2         # nucleus 2 zeeman
        )
H_Z = sp.diag(*entries) # Hamiltonian 
\end{lstlisting} 

This generates the final \verb|SymPy| 16$\times$16 Zeeman Hamiltonian. 

\newpage 
\section{Hyperfine Hamiltonian}

\subsection{Two-electron + One-nucleus}
To understand how to define the Hyperfine Hamiltonian for a two-electron $+$
two-nuclear system, we start with a simpler two-electron (a, b)  $+$ one-nucleus
(I) system. In this system, the Hyperfine Hamiltonian is defined as 
\begin{align*}
    \hat{H}_{HF} &= \hat{S}_a \cdot A_a \cdot \hat{I} + \hat{S}_b \cdot A_b \cdot
    \hat{I} \\ 
                 &= \sum_{k\in\{a, b\}}^{} A_{kx} S_{kx} I_x + A_{ky} S_{ky}I_y
                 + A_{kz}S_{kz}I_z. 
\end{align*}

Using ladder operator formalism, where 
\begin{align} \label{raisinglower}
    \hat{S}_x = \frac{1}{2} \left( \hat{S}_+ + \hat{S}_- \right) \qquad
    \hat{S}_y = \frac{1}{2i}\left( \hat{S}_+ - \hat{S}_- \right),  
\end{align} \vspace{3px}

and likewise for $\hat{I}_x, \hat{I}_y$ acting on the nuclear states, we have 

\begin{align*}
    \hat{H}_{HF} &= \sum_{k}^{} A_{kx} \left( \frac{1}{2} \left( S_{k+} + S_{k-}
    \right)   \right) \left( \frac{1}{2}\left( \hat{I}_+ + \hat{I}_- \right)
    \right)  \\ 
                 &\qquad+ A_{ky} \left( \frac{1}{2i} \left( S_{k+} - S_{k-} \right)
                 \right) \left( \frac{1}{2i} \left( I_+ - I_- \right)   \right)
                 + A_{kz} S_{kz}I_z \\ 
                 &= \sum_{k}^{} \frac{A_{kx}}{4} \left( S_{k+}I_+ + S_{k+}I_- +
                 S_{k-}I_+ + S_{k-}I_- \right) \\ 
                 &\qquad- \frac{A_{ky}}{4} \left( S_{k+}I_+ - S_{k+}I_- - S_{k-}I_+ +
                 S_{k-}I_- \right) + A_{kz} S_{kz} I_z.  
\end{align*}

And now since 

\[
\hat{S_\pm}|s, m\rangle = \hbar \sqrt{s(s+1) - m(m\pm 1)} |s, m\pm 1\rangle, 
\] \vspace{3px}

we can derive a full expression for $\hat{H}_{HF}|m_a, m_b, m_I\rangle$. Note
$s=1/2$ is implicit and we are working in the Zeeman $|\uparrow,
\downarrow\rangle$ basis to make things easier for now. \\ 

We calculate an expression for $H_{HF}$ that works for all 8 basis states in the
two-electron + one-nucleus system. On the individual electron states, we see 
\begin{align*}
  & S_+|\uparrow\rangle = 0 &&S_+ |\downarrow\rangle = \hbar|\uparrow\rangle \\ 
  & S_-|\downarrow\rangle = 0 &&S_- |\uparrow\rangle = \hbar|\downarrow\rangle
  \\ & &&S_z|m\rangle = m\hbar|m\rangle, \quad m\in\{\uparrow, \downarrow\} . 
\end{align*}

and likewise for the nuclear state, 
\begin{align*}
  & I_+|\Uparrow\rangle = 0 &&I_+ |\Downarrow\rangle = \hbar|\Uparrow\rangle \\ 
  & I_-|\Downarrow\rangle = 0 &&I_- |\Uparrow\rangle = \hbar|\Downarrow\rangle
  \\ & &&I_z|m_I\rangle = m_I\hbar|m_I\rangle, \quad m_I\in\{\Uparrow, \Downarrow\} . 
\end{align*}

Therefore, there are five nontrivial (that don't end up $=0$) operator products: 
\begin{align*}
    &\text{product} &&\text{nonzero for} &&&\text{result} \\  
    &S_{k+}I_+ &&|\downarrow, \Downarrow\rangle &&&\hbar^2|\uparrow,
    \Uparrow\rangle \\ 
    &S_{k+}I_- &&|\downarrow, \Uparrow\rangle &&&\hbar^2|\uparrow,
    \Downarrow\rangle \\ 
    &S_{k-}I_+ &&|\uparrow, \Downarrow\rangle &&&\hbar^2|\downarrow,
    \Uparrow\rangle \\ 
    &S_{k-}I_{-} &&|\uparrow, \Uparrow\rangle &&&\hbar^2|\downarrow,
    \Downarrow\rangle \\ 
    &S_{kz}I_z &&\text{any } m_{a,b}m_I &&& \hbar^2m_{a, b}m_I|m_{a,b}m_I\rangle  
\end{align*}

Hence, using this table, we can construct the Hamiltonian: 
\begin{align*}
    \hat{H}_{HF}|m_a, m_b, m_I\rangle &= \sum_{k\in\{a, b\}} \frac{\hbar^2}{4}\left( A_{kx} - A_{ky}
    \right) \hat{S}_{k+}\hat{I}_+ \delta_{m_k, m_I}|-m_k, \_, -m_I\rangle \\ 
                                      &\quad\qquad+ \frac{\hbar^2}{4}\left( A_{kx} - A_{ky}
                                      \right) \hat{S}_{k+}\hat{I}_-\delta_{m_k,
                                      -m_I}|-m_k, \_, -m_I\rangle \\ 
                                      &\quad\qquad + A_{kz}\hbar^2m_k m_I |m_k, \_,
                                      m_I\rangle. 
\end{align*} 

Expanding the summation, 
\begin{align*}
    \hat{H}_{HF}|m_a, m_b, m_I\rangle &= \hbar^2\left( A_{az}m_a m_I +
    A_{bz}m_bm_I \right) |m_a, m_b, m_I\rangle \\ 
                                      &+ \frac{\hbar^2}{4} \left[ \left( A_{ax}
                                          - A_{ay} \right) \delta_{m_a, m_I} +
                                          \left( A_{ax} + A_{ay} \right)
                                      \delta_{m_a, -m_I}\right] |-m_a, m_b,
                                      -m_I\rangle \\ 
                                      &+ \frac{\hbar^2}{4}\left[ \left( A_{bx} -
                                      A_{by} \right) \delta_{m_b, m_I} + \left(
                              A_{bx} + A_{by} \right) \delta_{m_b, -m_I} \right]
                              |m_a, -m_b, m_I\rangle. 
\end{align*}

where $m_a, m_b, m_I$ all run over $\{-\frac{1}{2}, \frac{1}{2}\}$. The equation
above defines the Hyperfine Hamiltonian in the Zeeman Basis for the full
8-basis-state two-electron + one-nuclear system. Evolving it to the two-electron + two-nuclear
system is straightforward. 

\subsection{Two-electron + Two-nuclei}

For two electrons $k\in\{a, b\}$ and two nuclei $p\in\{1, 2\}$, 

 \[
     \hat{H}_{HF} = \sum_{k\in\{a, b\}}^{} \sum_{p=1}^2 \left(
     A_{kpx}S_{kx}I_{px} + A_{kpy}S_{ky}I_{py} + A_{kpz}S_{kz}I_{pz} \right).   
\] \vspace{3px}

Following the same procedure as the simpler system, we get the full Hyperfine
Hamiltonian definition: 
\begin{equation} \label{hyperfine}
    \begin{split}
    \hat{H}_{HF} |&m_a, m_b, m_{I_1}, m_{I_2}\rangle  \\
                       &= \hbar^2 \sum_{p=1}^{2} \left(A_{apz}m_am_{I_p} +
                       A_{bpz}m_bm_{I_p} \right) |m_a, m_b, m_{I_1},
                       m_{I_2}\rangle  \\ 
                       &+\frac{\hbar^2}{4}\left[\left( A_{a1x}-A_{a1y} \right)
                           \delta_{m_a, m_{I_1}} + \left( A_{a1x} + A_{a1y}
                           \right) \delta_{m_a, -m_{I_1}}  \right] |-m_a, m_b,
                           -m_{I_1}, m_{I_2}\rangle \\ 
                       &+\frac{\hbar^2}{4}\left[\left( A_{a2x}-A_{a2y} \right)
                           \delta_{m_a, m_{I_2}} + \left( A_{a2x} + A_{a2y}
                           \right) \delta_{m_a, -m_{I_2}}  \right] |-m_a, m_b,
                            m_{I_1}, -m_{I_2}\rangle \\ 
                       &+\frac{\hbar^2}{4}\left[\left( A_{b1x}-A_{b1y} \right)
                           \delta_{m_b, m_{I_1}} + \left( A_{b1x} + A_{b1y}
                           \right) \delta_{m_b, -m_{I_1}}  \right] |m_a, -m_b,
                           -m_{I_1}, m_{I_2}\rangle \\ 
                       &+\frac{\hbar^2}{4}\left[\left( A_{b2x}-A_{b2y} \right)
                           \delta_{m_b, m_{I_2}} + \left( A_{b2x} + A_{b2y}
                           \right) \delta_{m_b, -m_{I_2}}  \right] |m_a, -m_b,
                           m_{I_1}, -m_{I_2}\rangle.   
\end{split}
\end{equation}


However, remember this is in the Zeeman $\uparrow, \downarrow$ basis. We need to
convert it to the coupled triplet-singlet basis we used in the Zeeman
Hamiltonian. Therefore, we apply a Clebsch-Gordan (CG) Unitary Transformation. \\ 

First we have 

 \[
\begin{pmatrix}
  |\uparrow \uparrow\rangle \\ |\uparrow\downarrow\rangle \\ |\downarrow
  \uparrow\rangle \\ |\downarrow\downarrow\rangle
\end{pmatrix} = \begin{pmatrix}
1 & 0 & 0 & 0 \\ 
0 & \frac{1}{\sqrt{2}} & \frac{1}{\sqrt{2}} & 0 \\ 
0 & \frac{1}{\sqrt{2}} & -\frac{1}{\sqrt{2}} & 0 \\ 
0 & 0 & 0 & 1 
\end{pmatrix}  \begin{pmatrix}
    |11\rangle \\ |10\rangle \\ |00\rangle \\ |1-1\rangle
\end{pmatrix},
\] \vspace{3px}

defining a transformation between the electron Zeeman $\to$ coupled basis.
Therefore, defining 

\[
 W = \begin{pmatrix}
     1 & 0 & 0 & 0 \\ 
     0 & \frac{1}{\sqrt{2}} & \frac{1}{\sqrt{2}} & 0 \\ 
     0 & \frac{1}{\sqrt{2}} & -\frac{1}{\sqrt{2}} & 0 \\ 
     0 & 0 & 0 & 1 
 \end{pmatrix} \otimes \begin{pmatrix}
     1 & 0 & 0 & 0 \\ 
     0 & 1 & 0 & 0 \\ 
     0 & 0 & 1 & 0 \\ 
     0 & 0 & 0 & 1 
 \end{pmatrix} 
\] \vspace{3px}

where we $\otimes I_4$ by the 4$\times$4 Identity because we leave the nuclear
$|m_{I_1}, m_{I_2}\rangle$ in the Zeeman basis and only couple the electrons.
Then, once we have the Hyperfine Hamiltonian in the \textit{Zeeman Basis}, we
can apply the transformation 
\begin{align}\label{cg} 
    \hat{H}_{HF_\text{coupled}} = W \left( \hat{H}_{HF_\text{zeeman} } \right)
    W^\dagger,  
\end{align} \vspace{3px}

to get the final Hyperfine Hamiltonian in the correct coupled basis. 

\subsection{Programming the Hamiltonian}

We first store Equation \ref{hyperfine} as a \verb|SymPy| 16$\times$16 matrix,
which is in the Zeeman $|m_a, m_b, m_{I_1}, m_{I_2}\rangle $ Basis. \\

We also have to ensure that after the CG transformation, shown in
Equation \ref{cg}, the now matrix coupled-basis ordering is the identical to the
coupled-basis ordering of the Zeeman Hamiltonian. Maintaining the order of the
coupled-basis is crucial. \\ 

The following code generates the 16 Zeeman basis states in correct order: 

\begin{lstlisting}[language=Python] 
def _generate_zeeman_basis(self): 
    # Builds the Zeeman |m_a, m_b, m_I1, m_I2> basis with m = +/-1/2.

    half = sp.Rational(1, 2)
    return [
        (m_a, m_b, m_I1, m_I2)
        for m_a in (half, -half)
        for m_b in (half, -half)
        for m_I1 in (half, -half)
        for m_I2 in (half, -half)
    ] 
... 
self.basis_ze = self._generate_zeeman_basis()
\end{lstlisting} 

\vspace{5px}
Now, to emulate Equation \ref{hyperfine}, we write the following functions:
\vspace{5px}

\begin{lstlisting}[language=Python] 
def _delta_parallel(self, m_e, m_I):
    return m_e == m_I

def _delta_antiparallel(self, m_e, m_I):
    return m_e == -m_I

def _action_on_ket(self, ket):
    # Calculates H_HF |ket> in Zeeman basis, returns dict {new_ket: coeff}.
    
    m_a, m_b, m_I1, m_I2 = ket
    psi = {}

    # diagonal S_z I_z term: A * m_e * m_I
    diag = (
        self.A['AA1Z'] * m_a * m_I1 +
        self.A['AB1Z'] * m_b * m_I1 +
        self.A['AA2Z'] * m_a * m_I2 +
        self.A['AB2Z'] * m_b * m_I2
    )
    psi[ket] = diag

    # off-diagonal flip-flop S_x I_x + S_y I_y -> (A_x - A_y)/4 etc.
    def add_flip(m_e, m_I, x_key, y_key, new_ket):
        if self._delta_parallel(m_e, m_I):
            coeff = (self.A[x_key] - self.A[y_key]) / 4
            psi[new_ket] = coeff
        elif self._delta_antiparallel(m_e, m_I):
            coeff = (self.A[x_key] + self.A[y_key]) / 4
            psi[new_ket] = coeff

    add_flip(m_a, m_I1, 'AA1X', 'AA1Y', (-m_a,  m_b, -m_I1,  m_I2))
    add_flip(m_a, m_I2, 'AA2X', 'AA2Y', (-m_a,  m_b,  m_I1, -m_I2))
    add_flip(m_b, m_I1, 'AB1X', 'AB1Y', ( m_a, -m_b, -m_I1,  m_I2))
    add_flip(m_b, m_I2, 'AB2X', 'AB2Y', ( m_a, -m_b,  m_I1, -m_I2))

    return psi 
\end{lstlisting} 

\verb|_delta_parallel| and \verb|_delta_antiparallel| are the delta functions
shown in Equation \ref{hyperfine} and allow us which of $\frac{1}{4}(A_{kpx} \pm
A_{kpy})$ are the coefficients of the superposed kets.  \\ 

We now perform the Hyperfine Hamiltonian (still Zeeman basis) on every Zeeman
basis-state, building the full Hyperfine Zeeman Matrix as we go. 

\begin{lstlisting}[language=Python] 
def _build_zeeman_matrix(self):
# Return the full 16x16 hyperfine template in the Zeeman basis.

size = len(self.basis_ze)
H = sp.MutableDenseMatrix(size, size, lambda *_: 0)
for j, ket in enumerate(self.basis_ze):
    action = self._action_on_ket(ket)
    for new_ket, coeff in action.items():
        i = self.index_ze[new_ket]
        H[i, j] = coeff
return H.as_immutable()
... 
self.H_ze = self._build_zeeman_matrix() 
\end{lstlisting} 

Now all that is left is to apply the CG transformation (Equation
\ref{cg}). We first build the $W$ via defining the 4$\times$4 CG matrix and taking the
tensor product with the 4$\times$4 Identity: 

\begin{lstlisting}[language=Python] 
def _build_cg_unitary(self):
    # Builds the unitary for two-electron Clebsch-Gordan transform otimes identity on nuclei.

    half = sp.sqrt(sp.Rational(1, 2))
    U = sp.Matrix([
        [1,    0,     0,    0],
        [0,  half,  half,   0],
        [0,  half, -half,   0],
        [0,    0,     0,    1]
    ])
    I4 = sp.eye(4)
    return sp.kronecker_product(U, I4)
... 
self.W = self._build_cg_unitary() 
\end{lstlisting} 

Finally, we get the Hyperfine Hamiltonian in the same coupled-basis as the
Zeeman Hamiltonian: 

\begin{lstlisting}[language=Python] 
self.H_coup = sp.simplify(self.W * self.H_ze * self.W.H)  # Hyperfine Hamiltonian 
\end{lstlisting} 

\subsection{Converting to Isotropic Hamiltonian}

We have the full, correct, Hyperfine Hamiltonian. However, for computational
feasibility, we will assume the hyperfine effect is \textit{isotropic}, that is, 
we will assume 
\begin{align*}
    A_{a1x} &= A_{a1y} = A_{a1} \\ 
    A_{b1x} &= A_{b1y} = A_{b1},
\end{align*}

and likewise for the second nuclear hyperfine interaction. Hence, we implement 

\begin{lstlisting}[language=Python] 
def convert_to_isotropic(H_sym: sp.Matrix) -> sp.Matrix:
    # Substitute anisotropic A_xyz -> isotropic Aa, Ab values.

    subs_map = {
        A_a1x: Aa1, A_a1y: Aa1, A_a1z: Aa1,
        A_a2x: Aa2, A_a2y: Aa2, A_a2z: Aa2,
        A_b1x: Ab1, A_b1y: Ab1, A_b1z: Ab1,
        A_b2x: Ab2, A_b2y: Ab2, A_b2z: Ab2,
    }
    return sp.simplify(H_sym.subs(subs_map)) 
\end{lstlisting} 

This returns the final isotropic Hyperfine Hamiltonian we will be using for
analysis. 

\newpage
\section{Zero-Field Splitting}

The Zero-Field Splitting (ZFS) Hamiltonian, also known as the Dipole-Dipole
Hamiltonian, is defined by
\begin{align} \label{zfs_start}
    \hat{H}_{ZFS} = D_1\left(\hat{S}_z^2 - \frac{1}{3}\left(\hat{S}_x^2 + \hat{S}_y^2 +
    \hat{S}_z^2 \right) \right) + D_2\left(\hat{S}_x^2 - \hat{S}_y^2\right).
\end{align} \vspace{3px}

Using raising/lowering operator formalism again, we have 
\begin{align*}
    \hat{S}_+|s, m\rangle &= \frac{1}{2}\left(\hat{S}_x + i \hat{S}_y\right) |s,
  m\rangle = \sqrt{s(s+1) - m(m+1)}|s, m+1\rangle \\ 
    \hat{S}_-|s, m\rangle &= \frac{1}{2}\left(\hat{S}_x - i \hat{S}_y\right) |s,
    m\rangle = \sqrt{s(s+1) - m(m-1)}|s, m-1\rangle.
\end{align*}

Hence, we retrieve Equation \ref{raisinglower} and find 
\begin{align*}
    \hat{S}_x &= \frac{1}{2}\left( \hat{S}_+ + \hat{S}_- \right) \\ 
    \hat{S}_y &= \frac{1}{2i}\left(\hat{S}_+ - \hat{S}_- \right) \\ 
    \hat{S}_x^2 + \hat{S}_y^2 = \hat{S}_x^2 - \hat{S}_y^2 &= \frac{1}{2}\left(
    \hat{S}_+^2 + \hat{S}_-^2 \right).  
\end{align*}

Substituting this into Equation \ref{zfs_start}, we get 
\begin{align} \label{zfs_state}
    \hat{H}_{ZFS} |s, m\rangle &= D_1m^2|s, m\rangle - \frac{D_1}{3}(s(s+1))|s,
    m\rangle \\ 
                               &+ \frac{D_2}{2}(s(s+1) - m(m+1))|s, m+2\rangle
                            \\ &+ \frac{D_2}{2}(s(s+1) - m(m-1))|s, m-2\rangle, 
\end{align}\vspace{3px} 

where the last two terms highlight the forbidden $m\pm 2$ transitions that give
rise to the half-field response. This Hamiltonian acts on the coupled basis,
akin to the Zeeman Hamiltonian. Hence, we translate the above ZFS
Hamiltonian definition as follows: 

\begin{lstlisting}[language=Python] 
electron_states = [(1, 1), (1, 0), (0, 0), (1, -1)]

# build 4x4 electron ZFS Hamiltonian 
H_elec = sp.zeros(4)
for i, (s, m) in enumerate(electron_states):
    # Diagonal term: D1 * m^2 - (D1/3) * s(s+1)
    H_elec[i, i] = D1*m**2 - (D1/3)*s*(s+1)
    # Off-diagonal terms coupling m -> m+/-2
    for dm, expr in [(2, D2/2*(s*(s+1) - m*(m+1))),
                     (-2, D2/2*(s*(s+1) - m*(m-1)))]:
        newm = m + dm
        if (s, newm) in electron_states:
            j = electron_states.index((s, newm))
            H_elec[i, j] = expr

# 4x4 identity nuclear ZFS Hamiltonian 
I_nuc = sp.eye(4)

# Final ZFS Hamiltonian 
H_ZFS = sp.kronecker_product(H_elec, I_nuc)
\end{lstlisting} 

where we again take the tensor product $H_\text{electron ZFS} \otimes I_4$ to preserve basis ordering. This gives the final ZFS/Dipole-Dipole
Hamiltonian in the correct coupled-basis. 

\newpage
\section{Exchange Hamiltonian}

The isotropic Exchange Hamiltonian between electrons is defined by 

\[
    \hat{H}_{EX} = -J \hat{S}_a \cdot \hat{S}_b,
\] \vspace{3px}

where we add a $-1$ factor to yield $-\frac{J}{4}$ for triplet states and
$+\frac{3J}{4}$ for singlet states. Therefore, its action on a coupled $|s, m\rangle$ state of two electrons is
given by 

\[
    \hat{H}_{EX} |s, m\rangle = -J \left( \hat{S}_a \cdot \hat{S}_b \right) |s,
    m\rangle = -J \left( \frac{s(s+1) - \frac{3}{2}}{2} \right) |s, m\rangle.  
\] \vspace{3px}

Because this already acts on the coupled-basis (nuclear $|m_I\rangle $are left
alone), programming this Hamiltonian is the easiest. 

\begin{lstlisting}[language=Python] 
# same electron basis ordering
electron_states = [(1, 1), (1, 0), (0, 0), (1, -1)]

# build 4x4 electron exchange Hamiltonian 
H_ex_elec = sp.zeros(4)
for i, (s, m) in enumerate(electron_states):
    H_ex_elec[i, i] = -J * (s * (s + 1) - 1.5) / 2

# 4x4 nuclear identity 
I_nuc = sp.eye(4)

# Exchange Hamiltonian 
H_EX = sp.kronecker_product(H_ex_elec, I_nuc)
\end{lstlisting} 

Here, to generate the full 16$\times$16 Exchange Hamiltonian that acts on the
full two-electron $+$ two-nuclei system, we again take the tensor product of the
4$\times$4 electron Exchange Hamiltonian with the  4$\times$4 nuclear Identity
matrix to yield the full Exchange Hamiltonian on the system. 

\vfill 
The Full 4H-Silicon Carbide Spin Hamiltonian is finally, 

\begin{lstlisting}[language=Python] 
H_SPIN = H_zeeman + H_hyperfine + H_zfs + H_exchange.
\end{lstlisting}

I.e., the sum of all previously defined Hamiltonians defined in same
coupled-basis throughout. 
\vfill

\section{Basis Appendix} \label{Basis} 

\subsection{Zeeman Basis}

The Zeeman-Basis used for calculating the Hyperfine Hamiltonian is defined by 
\begin{align*}
  &0 &&\bigl|+\tfrac{1}{2}, +\tfrac{1}{2}, +\tfrac{1}{2}, +\tfrac{1}{2}\bigr\rangle\\
  &1 &&\bigl|+\tfrac{1}{2}, +\tfrac{1}{2}, +\tfrac{1}{2}, -\tfrac{1}{2}\bigr\rangle\\
  &2 &&\bigl|+\tfrac{1}{2}, +\tfrac{1}{2}, -\tfrac{1}{2}, +\tfrac{1}{2}\bigr\rangle\\
  &3 &&\bigl|+\tfrac{1}{2}, +\tfrac{1}{2}, -\tfrac{1}{2}, -\tfrac{1}{2}\bigr\rangle\\
  &4 &&\bigl|+\tfrac{1}{2}, -\tfrac{1}{2}, +\tfrac{1}{2}, +\tfrac{1}{2}\bigr\rangle\\
  &5 &&\bigl|+\tfrac{1}{2}, -\tfrac{1}{2}, +\tfrac{1}{2}, -\tfrac{1}{2}\bigr\rangle\\
  &6 &&\bigl|+\tfrac{1}{2}, -\tfrac{1}{2}, -\tfrac{1}{2}, +\tfrac{1}{2}\bigr\rangle\\
  &7 &&\bigl|+\tfrac{1}{2}, -\tfrac{1}{2}, -\tfrac{1}{2}, -\tfrac{1}{2}\bigr\rangle\\
  &8 &&\bigl|-\tfrac{1}{2}, +\tfrac{1}{2}, +\tfrac{1}{2}, +\tfrac{1}{2}\bigr\rangle\\
  &9 &&\bigl|-\tfrac{1}{2}, +\tfrac{1}{2}, +\tfrac{1}{2}, -\tfrac{1}{2}\bigr\rangle\\
  &10 &&\bigl|-\tfrac{1}{2}, +\tfrac{1}{2}, -\tfrac{1}{2}, +\tfrac{1}{2}\bigr\rangle\\
  &11 &&\bigl|-\tfrac{1}{2}, +\tfrac{1}{2}, -\tfrac{1}{2}, -\tfrac{1}{2}\bigr\rangle\\
  &12 &&\bigl|-\tfrac{1}{2}, -\tfrac{1}{2}, +\tfrac{1}{2}, +\tfrac{1}{2}\bigr\rangle\\
  &13 &&\bigl|-\tfrac{1}{2}, -\tfrac{1}{2}, +\tfrac{1}{2}, -\tfrac{1}{2}\bigr\rangle\\
  &14 &&\bigl|-\tfrac{1}{2}, -\tfrac{1}{2}, -\tfrac{1}{2}, +\tfrac{1}{2}\bigr\rangle\\
  &15 &&\bigl|-\tfrac{1}{2}, -\tfrac{1}{2}, -\tfrac{1}{2}, -\tfrac{1}{2}\bigr\rangle
\end{align*}

where $|+\frac{1}{2}\rangle = |\uparrow\rangle$ and $-\frac{1}{2} =
|\downarrow\rangle$ (electron) or $|\Uparrow, \Downarrow\rangle$ (nuclear) \\

\vspace{5px} 
Indices 0-15 represent matrix left-to-right or top-to-bottom basis states. 


\newpage 

\subsection{Coupled Basis}

The Coupled-Basis used for calculating all individual Hamiltonians and summing
them together to generate the final Spin Hamiltonian is defined by 

\[
\begin{alignedat}{2}
 0\qquad\qquad\qquad\qquad & |1,1\rangle   \quad & \otimes \quad & |+\tfrac12,+\tfrac12\rangle\\
 1\qquad\qquad\qquad\qquad & |1,1\rangle   \quad & \otimes \quad & |+\tfrac12,-\tfrac12\rangle\\
 2\qquad\qquad\qquad\qquad & |1,1\rangle   \quad & \otimes \quad & |-\tfrac12,+\tfrac12\rangle\\
 3\qquad\qquad\qquad\qquad & |1,1\rangle   \quad & \otimes \quad & |-\tfrac12,-\tfrac12\rangle\\[4pt]
 4\qquad\qquad\qquad\qquad & |1,0\rangle   \quad & \otimes \quad & |+\tfrac12,+\tfrac12\rangle\\
 5\qquad\qquad\qquad\qquad & |1,0\rangle   \quad & \otimes \quad & |+\tfrac12,-\tfrac12\rangle\\
 6\qquad\qquad\qquad\qquad & |1,0\rangle   \quad & \otimes \quad & |-\tfrac12,+\tfrac12\rangle\\
 7\qquad\qquad\qquad\qquad & |1,0\rangle   \quad & \otimes \quad & |-\tfrac12,-\tfrac12\rangle\\[4pt]
 8\qquad\qquad\qquad\qquad & |0,0\rangle   \quad & \otimes \quad & |+\tfrac12,+\tfrac12\rangle\\
 9\qquad\qquad\qquad\qquad & |0,0\rangle   \quad & \otimes \quad & |+\tfrac12,-\tfrac12\rangle\\
10\qquad\qquad\qquad\qquad & |0,0\rangle   \quad & \otimes \quad & |-\tfrac12,+\tfrac12\rangle\\
11\qquad\qquad\qquad\qquad & |0,0\rangle   \quad & \otimes \quad & |-\tfrac12,-\tfrac12\rangle\\[4pt]
12\qquad\qquad\qquad\qquad & |1,-1\rangle  \quad & \otimes \quad & |+\tfrac12,+\tfrac12\rangle\\
13\qquad\qquad\qquad\qquad & |1,-1\rangle  \quad & \otimes \quad & |+\tfrac12,-\tfrac12\rangle\\
14\qquad\qquad\qquad\qquad & |1,-1\rangle  \quad & \otimes \quad & |-\tfrac12,+\tfrac12\rangle\\
15\qquad\qquad\qquad\qquad & |1,-1\rangle  \quad & \otimes \quad & |-\tfrac12,-\tfrac12\rangle
\end{alignedat}
\]

where $|+\frac{1}{2}\rangle = |\Uparrow\rangle$ and $-\frac{1}{2} =
|\Downarrow\rangle$ (nuclear-only).\\

\vspace{5px} 
Indices 0-15 represent matrix left-to-right or top-to-bottom basis states. 













\end{document}
